% Generated draft: requires mathematical and source review.
\documentclass[11pt]{article}
\usepackage[margin=25mm]{geometry}
\usepackage{fontspec}
\setmainfont{Latin Modern Roman}
\usepackage{amsmath,amssymb,mathtools,unicode-math}
\setmathfont{Latin Modern Math}
\usepackage{enumitem,needspace,xurl,hyperref}
\setlength{\parindent}{0pt}
\setlength{\parskip}{5pt}
\setlength{\emergencystretch}{3em}
\newcommand{\N}{\mathbb{N}}
\newcommand{\Z}{\mathbb{Z}}
\newcommand{\Q}{\mathbb{Q}}
\newcommand{\R}{\mathbb{R}}
\newcommand{\C}{\mathbb{C}}
\newcommand{\F}{\mathbb{F}}
\newcommand{\A}{\mathbb{A}}
\newcommand{\PP}{\mathbb{P}}
\newcommand{\E}{\mathbb{E}}
\newcommand{\eps}{\varepsilon}
\newcommand{\dd}{\,\mathrm{d}}
\begin{document}
\section*{UnsolvedMath ID: 701055}
\textbf{Source problem code:} catalogue entry
\subsection{Definitions and mathematical statement}
For a group \(G\), use the commutator convention \([x,y]=x^{-1}y^{-1}xy\), and define left-normed iterated commutators by
\[
[x,{}_1y]=[x,y],\qquad [x,{}_{k+1}y]=[[x,{}_ky],y]\quad(k\geq 1).
\]
For a positive integer \(n\), the group \(G\) is \(n\)-Engel if
\[
\forall x,y\in G,\qquad [x,{}_ny]=1.
\]
Define its lower central series by \(\gamma_1(G)=G\) and
\[
\gamma_{k+1}(G)=\langle [a,b]\mid a\in\gamma_k(G),\ b\in G\rangle\quad(k\geq 1).
\]
The group is nilpotent if \(\gamma_{c+1}(G)=\{1\}\) for some integer \(c\geq 0\), and finitely generated if \(G=\langle S\rangle\) for some finite subset \(S\subseteq G\).

The existence question of B. I. Plotkin asks whether there are a positive integer \(n\) and a finitely generated group \(G\) such that
\[
\bigl(\forall x,y\in G,\ [x,{}_ny]=1\bigr)
\quad\text{and}\quad
\bigl(\forall c\geq 0,\ \gamma_{c+1}(G)\neq\{1\}\bigr).
\]
Here \(n\) is a single finite bound valid for all pairs of elements. The source does not specify a particular value of \(n\).

The supplied comment dated 2017 states that finitely generated \(n\)-Engel groups are nilpotent for \(n=2,3,4\). It also reports that E. Rips gives a course of lectures devoted to constructing non-nilpotent finitely generated \(n\)-Engel groups for sufficiently large \(n\), and that the lectures are available on the Internet. No lecture reference, link, or numerical threshold is supplied. These statements are retained as dated source commentary, without independent verification or an inference that the requested construction has been established. The local record labels the problem open; that label is supplied registry metadata.
\subsection{Short English statement}
Does there exist a finitely generated, non-nilpotent group satisfying the \(n\)-Engel identity for some fixed positive integer \(n\)? The supplied 2017 comment records nilpotence for \(n=2,3,4\) and reports lectures by E. Rips concerning constructions for sufficiently large \(n\); the lecture reference and threshold are unspecified.
\subsection{Sources}
\begin{itemize}
\item Local UnsolvedMath registry, permanent ID \(701055\), source code \texttt{catalogue entry}, title \emph{Kourovka Notebook Problem 14.70}; supplied source record, updated 2026-10-16, upstream revision \texttt{d371ee917ac9647fbb06cd6fc529deec6f3dda82}.
\item B. I. Plotkin, problem 14.70, \emph{The Kourovka Notebook (Unsolved Problems in Group Theory)}, 14th issue (1999). The supplied record identifies the containing edition as the 21st issue edition (2026), \texttt{arXiv:1401.0300v48}. The collection is attributed to editors E. I. Khukhro and V. D. Mazurov. These bibliographic details are supplied attribution and have not been independently verified.
\item Supplied comment dated 2017: nilpotence for finitely generated \(n\)-Engel groups when \(n=2,3,4\), and a report of Internet lectures by E. Rips devoted to constructions for sufficiently large \(n\). No separate citation or URL for the lectures or the small-\(n\) results is supplied.
\end{itemize}
\end{document}
